\subsection{Coulomb-WW als Beispiel}
\stp{Heisenberg-Bild:}
\begin{equation}
  H_{\mathrm{int}}^{H}=\frac{1}{2}\int d^{3}r\int d^{3}r\mkern1mu'\,\psi_{H}^{\dagger}(\vec r,t)\,\psi_{H}^{\dagger}(\vec r\mkern3mu',t)\,V(\vec r-\vec r\mkern3mu')\,\psi_{H}(\vec r\mkern3mu',t)\,\psi_{H}(\vec r,t)
\end{equation}

\stp{Wechselwirkungsbild:}
\begin{equation}
  \begin{aligned}
    \psi_{H}(\vec r,t)           & =S^{\dagger}(t,t_{0})\,\psi_{0}(\vec r,t)\,S(t,t_{0})           \\
    \psi_{H}^{\dagger}(\vec r,t) & =S^{\dagger}(t,t_{0})\,\psi_{0}^{\dagger}(\vec r,t)\,S(t,t_{0})
  \end{aligned}
\end{equation}
Damit folgt
\begin{equation}
  \begin{aligned}
    H_{\mathrm{int}}^{H}(t)   & =S^{\dagger}(t,t_{0})\,H_{\mathrm{int}}^{0}(t)\,S(t,t_{0})                                                                                                                                       \\
    H_{\mathrm{int}}^{(0)}(t) & =\frac{1}{2}\int d^{3}r\int d^{3}r\mkern1mu'\,\psi_{0}^{\dagger}(\vec r,t)\,\psi_{0}^{\dagger}(\vec r\mkern3mu',t)\,V(\vec r-\vec r\mkern3mu')\,\psi_{0}(\vec r\mkern3mu',t)\,\psi_{0}(\vec r,t)
  \end{aligned}
\end{equation}
Fuer das Coulomb-Potential verwenden wir
\begin{equation}
  V(\vec r-\vec r\mkern3mu')=\frac{e^{2}}{4\pi\varepsilon_{0}}\frac{1}{\lvert\vec r-\vec r\mkern3mu'\rvert}
\end{equation}
und in kompakter Schreibweise
\begin{equation}
  V(1-1\mkern1mu')=V(\vec r_{1}-\vec r_{1\mkern3mu'})\,\delta(t_{1}-t_{1\mkern1mu'})
\end{equation}

Fuer die Stoerungsentwicklung von $S$ wird benoetigt:
\begin{equation}
  \int dt\,H_{\mathrm{int}}^{0}(t)=\frac{1}{2}\int d1\int d1\mkern1mu'\,\psi_{0}^{\dagger}(1)\,\psi_{0}^{\dagger}(1\mkern1mu')\,V(1-1\mkern1mu')\,\psi_{0}(1\mkern1mu')\,\psi_{0}(1)
\end{equation}
Entwicklung von $G(1,1\mkern3mu')$ nach Potenzen der WW:
\begin{equation}\label{f:entw_G(1,1')}
  \begin{aligned}
    G(1,1\mkern1mu')       & =G^{(0)}(1,1\mkern1mu')+G^{(1)}(1,1\mkern1mu')+G^{(2)}(1,1\mkern1mu')+\ldots                                                                                              \\
    G^{(0)}(1,1\mkern1mu') & =-\frac{i}{\hbar}\frac{1}{\erw{S}}\erw{\mathcal T\left\{\psi_{0}(1)\psi_{0}^{\dagger}(1\mkern1mu')\right\}}                                                               \\
    G^{(1)}(1,1\mkern1mu') & =-\frac{i}{\hbar}\frac{1}{\erw{S}}\int d\bar t_{1}\,\erw{\mathcal T\left\{\psi_{0}(1)\psi_{0}^{\dagger}(1\mkern1mu')H_{\mathrm{int}}^{0}(\bar t_{1})\right\}}             \\
                           & =\left(-\frac{i}{\hbar}\right)^{2}\frac{1}{\erw{S}}\frac{1}{2}\int d2\int d2\mkern1mu'\,V(2-2\mkern1mu')                                                                  \\
                           & \quad\times\erw{\mathcal T\left\{\psi_{0}(1)\psi_{0}^{\dagger}(1\mkern1mu')\psi_{0}^{\dagger}(2)\psi_{0}^{\dagger}(2\mkern1mu')\psi_{0}(2\mkern1mu')\psi_{0}(2)\right\}}.
  \end{aligned}
\end{equation}

Das sind also viele Feldoperatoren, die wir mit Hilfe von Wick's Theorem in Produkte von $G^{(0)}$-Funktion umschreiben koennen. Zudem ist anzumerken, dass die Menge an Thermen immer größer wird, für ein genaues Ergebnis allerdings unendlich viele Terme notwendig sind.